\section{A certified class bound for the missing e133 reflection field}
\label{sec:cubic-e133-class-certificate}
The certificate below replaces unsuccessful class-quotient certification
by exact unit saturation and a rigorously bounded analytic interval.
No provisional class-group structure or complete unit basis is assumed.

\begin{theorem}[AC-G131-e133-class: the missing sextic class sector]
\label{AC-G131-e133-class}
Let $L=\mathbb Q(y)$ with $y^6=409115375$.
Then $\operatorname{Cl}(L)[269]=0$.
For $M=L(\sqrt{-3})$, the nontrivial cubic-character part of
$\operatorname{Cl}(M)[269]$ also vanishes.
This supplies a finite class-sector certificate; it does not, by itself,
construct a new e133 native host, modular support or model bound.
\end{theorem}
\begin{proof}
Put $p=269$, $d=5\cdot23\cdot269=30935$, and
\[
 L=\mathbb Q(y),\quad y^6=409115375=115^3\cdot269,
 \quad t=y^2/115,\quad s=y^3/115.
\]
Then $t^3=269$, $s^2=d$, and
$y=s/t$, so $L=\mathbb Q(t,\sqrt d)$.
The displayed polynomial is Eisenstein at269.  Its signature is $(2,2)$,
and it has only the roots of unity $\{1,-1\}$ because it has a real embedding.

\proofstep{1. An exact saturated unit lattice.}
Let $u_1,u_2,u_3$ be the three rational polynomials in $y$ in
\texttt{saturation\_certificate.json}.  Their complete coefficients and
monic annihilating polynomials are in that file.  For orientation,
\[
 u_1=155481+\frac{884}{115}y^3,
 \qquad
 u_2=-270+\frac{739}{115}y^2-\frac{108}{13225}y^4.
\]
Their respective annihilating polynomials are
$x^2-310962x+1$ and $x^3+810x^2+64626984x+1$.
The third unit has the recorded monic reciprocal degree-six polynomial
with constant1.  The independent verifier evaluates all three polynomials
in $\mathbb Q[y]/(y^6-409115375)$ and obtains zero exactly.
Thus each $u_i$ is an algebraic integer with integral inverse.
No assertion that PARI's units are a full unit basis is used.

There are three residue homomorphisms with the following data:
\[
\begin{array}{c|c|c|ccc}
\ell&y\bmod\ell&\zeta\bmod\ell&d_1&d_2&d_3\\\hline
2153&451&102&169&226&225\\
2153&1702&102&100&226&44\\
3767&48&3739&42&87&263
\end{array}
\]
For each row, trial division verifies primality, $y$ is a simple root of
the defining polynomial modulo$\ell$, and all displayed unit denominators
are invertible.  Hensel's lemma gives an embedding into $\mathbb Q_\ell$.
The integer $\zeta$ has exact order269 and the exact modular identities are
\[
        u_i(y)^{(\ell-1)/269}=\zeta^{d_i}\pmod\ell.
\]
The verifier checks these identities for all30 supplied rows, including
these three.  The displayed $3\times3$ logarithm matrix has determinant
$99$ modulo269.  In particular its rows give an injection of the generated
unit lattice modulo269 into a product of three residue groups.

Write $U=\langle-1,u_1,u_2,u_3\rangle$.
The same matrix shows that $u_1,u_2,u_3$ are multiplicatively independent:
a primitive integral relation, including a possible sign, would give a
nonzero vector modulo269 in the matrix kernel.  Since the unit rank is3,
$U$ has finite index $m$ in $\mathcal O_L^\times$.
If269 divided $m$, there would be $v\notin U$ with
$v^{269}=(-1)^a\prod u_i^{b_i}$.
Every residue character kills269th powers and $-1$, so269 divides every
$b_i$.  Dividing by the corresponding product and sign gives a269th root
of unity in $L$.  It is1, contradicting $v\notin U$.  Therefore
\[
                     269\nmid m.
\]
This is an actual exact unit-saturation conclusion.  It does not yet say
anything about the class group.

\proofstep{2. Local fields, discriminant, and analytic factors.}
Put $K=\mathbb Q(\sqrt{-3})$ and let $\chi$ be either nontrivial cubic
character of $K(\sqrt[3]{269})/K$.
The normal closure of $\mathbb Q(t)$ has group$S_3$.
It meets the real quadratic field $\mathbb Q(\sqrt d)$ trivially,
since its sole quadratic subfield is$K$.
Let $\delta=\chi_\Delta$, where
\[
                  \Delta=4d=123740.
\]
The permutation representation for $L$ is
$(1\oplus\operatorname{Ind}_K^{\mathbb Q}\chi)\otimes(1\oplus\delta)$.
Consequently its Euler factors give
\[
 \zeta_L(s)=\zeta(s)L(s,\delta)L(s,\rho)L(s,\rho\delta),
 \qquad \rho=\operatorname{Ind}_K^{\mathbb Q}\chi.
\]
Here the two degree-two functions are the finite-order Hecke functions
of $\chi$ and $\chi\delta|_{G_K}$, respectively.

The number269 is a cube in $\mathbb Q_3$.  Indeed, writing
$-(1+3z)$ for a cube root reduces the equation to
$z+3z^2+3z^3=-30$, which has a solution in$\mathbb Z_3$ by
Hensel's lemma, its derivative being1 modulo3.
Thus the cubic character is trivial at3.  At269 it has tame order3,
and it is unramified at every other finite place of$K$.
The first degree-two conductor is
\[
                    Q_0=3\cdot269^2=217083.
\]
The twist has order6 tame inertia at269; it has scalar quadratic inertia
at5 and23, and the conductor-two quadratic scalar at2.
At3 it still has inertia $1\oplus\eta_K$.
Therefore its conductor is
\[
                  Q_1=2^4\cdot3\cdot5^2\cdot23^2\cdot269^2
                      =3\Delta^2=45934762800.
\]
These conductor exponents follow directly from the Artin-conductor
formula; at2 an unramified two-dimensional representation is twisted by
a scalar of conductor exponent2.

There is also a direct check of the field discriminant, independent of
any uncertified class calculation.  At269 the degree-six Eisenstein
extension is totally and tamely ramified, with different exponent5.
At2,5,23 the cubic algebra is unramified; the quadratic discriminant
contributes three times its exponents2,1,1.
At3 the cubic algebra is $\mathbb Q_3\times K_3$ and the quadratic field
is unramified of degree2, since $d$ is nonsquare modulo3.
The total exponent is therefore2.  All other places are unramified.
Hence
\[
 D_L=2^6\cdot3^2\cdot5^3\cdot23^3\cdot269^5
     =\Delta Q_0Q_1=1233892727411780376000.
\]

The coefficients $a_n$ of $L(s,\rho)$ are determined multiplicatively.
At269 they vanish in every positive prime-power degree, at3 they are1;
for other primes $r\equiv2\pmod3$ they are1 in even degree and0 in odd degree.
For other primes $r\equiv1\pmod3$, the trace is2 if269 is a cube modulo$r$ and$-1$
otherwise, and the determinant is1, giving the usual degree-two recurrence.
Thus $|a_n|\le d(n)\le n$.
The coefficients for the twisted function are exactly $a_n\delta(n)$.
This includes the bad primes: at2,5,23,269 the twisted factor is1,
and at3 it is $(1+3^{-s})^{-1}$ because$\delta(3)=-1$.

\proofstep{3. Primitive finite Fourier transforms, including the modulus four.}
The theta identity needed for the computation can be proved directly.
The imaginary quadratic ring $\mathcal O_K=\mathbb Z[\omega]$ has
class number one and unit group$\mu_6$: the imaginary-quadratic
Minkowski bound $(2/\pi)\sqrt3<2$ gives the class assertion.
The two Hecke characters above are primitive finite characters $\nu_f$
modulo$f\mathcal O_K$, with$f=269$ and$f=\Delta$, respectively.
They are trivial on$\mu_6$, as also follows from their being ideal
characters.  For the cubic factor the local character is trivial on
$\mu_3$ since$9\mid269+1$; norm characters are trivial on all$\mu_6$.

Here is the only new prime-power conductor check.  The local field
$K_2/\mathbb Q_2$ is unramified quadratic.  For $a\in\mathcal O_{K_2}$,
\[
 N_{K_2/\mathbb Q_2}(1+2a)\equiv1+2\operatorname{Tr}_{\mathbb F_4/\mathbb F_2}
                       (\bar a)\pmod4.
\]
The residue trace is onto, so the norm pullback of the conductor-four
quadratic character is nontrivial on$1+2\mathcal O_{K_2}$ and trivial
on$1+4\mathcal O_{K_2}$.  Thus its conductor is exactly$(4)$.
At5 and23 the residue norm is onto and the pulled-back quadratic character
is nontrivial; at269 its order2 combines with order3 to give order6.
This proves primitivity at every prime dividing$\Delta$.

Extend$\nu_f$ by zero on nonunits of$A=\mathcal O_K/f\mathcal O_K$.
The additive trace pairing is perfect since$(f,3)=1$.
At a unit frequency, substituting the inverse frequency in the finite
Fourier sum shows that the transform is a fixed Gauss constant times
the conjugate multiplicative character (up to the fixed conjugation
in the Euclidean pairing).
At a nonunit frequency, select a prime$\mathfrak p^a$ of the modulus
at which the frequency is nonunit.  If$a>1$, primitivity supplies a unit
$u\in1+\mathfrak p^{a-1}$ on which the character is nontrivial.
Multiplication by$u$ fixes that additive character because the frequency
has positive valuation.  If$a=1$, use any unit with nontrivial character;
the local frequency is zero.  By the Chinese remainder theorem choose
$u=1$ at every other prime.  Substitution by$u$ makes the Fourier sum
its product with a nontrivial scalar, so the sum is zero.
This proves the nonunit cancellation also at$2^2$, without assuming
a squarefree modulus.  Plancherel now gives Gauss magnitude
$\sqrt{|A|}=f$: there are$|A^\times|$ equal-magnitude nonzero
Fourier values and the norm square of the original function is
$|A^\times|$.

Put$c=2\pi/(\sqrt3 f)=2\pi/\sqrt Q$ and
\[
 \Theta(t)=\frac16\sum_{\alpha\in\mathcal O_K}
               \nu_f(\alpha)e^{-c|\alpha|^2t}
          =\sum_{n\ge1}b_ne^{-cnt}.
\]
The equality groups the six generators of each principal ideal.
Poisson summation on$f\mathcal O_K$ gives
\[
              \Theta(t)=\epsilon t^{-1}\overline{\Theta(1/t)},
                         \qquad |\epsilon|=1.
\]
Indeed the lattice area is$f^2\sqrt3/2$ and its dual is
$2/(f\sqrt{-3})\mathcal O_K$.  The two-dimensional Gaussian transform
is$\pi/(ct)$ times$\exp(-\pi^2|w|^2/(ct))$.
These constants change$ct$ to$c/t$ and give factor$1/(ft)$ before
the finite Gauss constant of magnitude$f$.
A fixed rotation or conjugation of the summation lattice preserves norms
and changes only$\epsilon$; the conjugate character gives the displayed
conjugate theta series.  This derives the exact normalization rather than
importing an unspecified functional-equation phase.
Both coefficient sequences here are real.
The interval computation proves $\Theta(1)\ne0$, hence$\epsilon=1$.
The nonprincipal finite character has no constant term, and the transformed
series gives exponential decay at zero as well as infinity.  Its Mellin
integral is therefore entire; the termwise integration below is absolute.
Splitting the Mellin integral at1 gives the exact identity
\[
 L(1)=\sum_{n\ge1}\frac{b_ne^{-cn}}n
       +c\sum_{n\ge1}b_n E_1(cn),
 \qquad E_1(x)=\int_x^\infty e^{-u}\frac{du}{u}.
\]
For a truncation after$N$, put$z=e^{-c}$.
Using $E_1(x)\le e^{-x}/x$ and$|b_n|\le n$, absolute bounds for the
three tails of the first sum, second sum and theta value are
\[
 \frac{z^{N+1}}{1-z},\qquad
 \frac{z^{N+1}}{c(1-z)},\qquad
 \frac{z^{N+1}((N+1)-Nz)}{(1-z)^2}.
\]
They are added as symmetric error balls; no alternating-series assumption
or unproved coefficient cancellation is used.

For the primitive even quadratic character the exact finite formula is
\[
 L(1,\delta)=-\frac1{\sqrt\Delta}
       \sum_{a=1}^{\Delta}\delta(a)
                    \log(2\sin(\pi a/\Delta)).
\]
The terms with$\delta(a)=0$ are omitted.  For completeness, $\delta$ is the primitive even product of the
quadratic characters of conductors$4,5,23,269$.
The finite Fourier cancellation just proved in one dimension gives its
Gauss norm$\sqrt\Delta$; evenness makes the Gauss sum real.
The separate rigorous finite cosine sum in
\texttt{verify\_quadratic\_gauss.json} is positive, with value
$351.766968318516\ldots$, so its sign is the positive square root.
The identity follows by expanding
$-\log|1-re^{2\pi ix}|$ for$0<r<1$, taking the finite character sum,
and then letting$r$ tend to1.  Bounded partial sums of the nonprincipal
character justify the Abel limit on the Dirichlet side; every surviving
logarithm has nonzero argument.  Thus no unstated Gauss sign convention
enters the formula.
The function is evaluated by pairing$a$ with$\Delta-a$.
The logarithmic regulator $R_U$ is evaluated from the two real embeddings
and one representative of each of the two complex pairs, with weights
1,1,2,2.  Omitting any one row gives its determinant.
Finally the analytic class-number formula gives
\[
 A=h_LR_L=\frac{\sqrt{D_L}}{8\pi^2}
               L(1,\delta)L(1,\rho)L(1,\rho\delta),
 \qquad \frac A{R_U}=\frac{h_L}{m}.
\]

\proofstep{4. An elementary denominator bound and exact isolation.}
This argument uses no conjectural class bound and no Friedman constant.
For a degree-$n$ field set
$P_n=\prod_{j=1}^n(4\binom nj+1)$.
If a unit has all conjugates at most
$\exp(\log2/[n(P_n+1)])$ in modulus, the characteristic polynomials
of its first$P_n+1$ powers have coefficients bounded by$2\binom nj$.
There are at most$P_n$ such integer polynomials, so two coincide.
Comparison of the largest root moduli forces that modulus to be1.
The norm forces every conjugate modulus to be1; the same bounded-polynomial
argument applied to all powers then proves torsion.
A nontorsion unit therefore has a logarithmic coordinate exceeding
$\log2/[n(P_n+1)]$.

The weighted logarithmic lattice has coordinate sum zero.  After omitting
one of its$r+1$ coordinates, a nonzero lattice vector consequently has
supremum norm greater than$\log2/[nr(P_n+1)]$.
Translates of an open cube of that side length are disjoint modulo the
lattice.  Comparing their volume with the covolume proves
\[
                    R_L>\frac1{[2nr(P_n+1)]^r}.
\]
For$n=6,r=3$, this gives
\[
 P_6=941878125,\qquad
 R_L> [36(P_6+1)]^{-3}.
\]
The rigorous regulator ball in the script has$476282<R_U<476283$.
Thus the integer$m=R_U/R_L$ is less than
\[
                      M=476283[36(P_6+1)]^3.
\]
If the full reviewed interval computation establishes
\[
                     \left|\frac A{R_U}-1008\right|<\frac1M,
\]
then the integer$h_L-1008m$ must be zero: otherwise its quotient by$m$
has modulus at least$1/m>1/M$.
Therefore$h_L=1008m$; the exact saturation in section1 and$269\nmid1008$
then imply$\mathrm{Cl}(L)[269]=0$.
This does not assert$h_L=1008$ or that the candidate units are complete.
The completed ball calculation uses900-bit precision for the regulator,
256-bit precision for the three $L$-values, $N=12000$ for$Q_0$ and
$N=5000000$ for$Q_1$.  Every tail above is added with outward rounding.
Its exact strict ball comparison proves
\[
 \left|A/R_U-1008\right|<1.2\cdot10^{-41},
 \qquad 1/M>5.38\cdot10^{-38}.
\]
The regulator determinant uses the actual exact polynomial units and
weights$1,1,2,2$; all four column sums contain zero.
The preserved output records every log interval, conductor, theta value,
$L$-value and absolute tail.  In particular its root-number tests give
$\Theta_0(1)=5.01709\ldots$ and
$\Theta_1(1)=-2.85482\ldots$, both bounded away from zero.
The assertions are comparisons of interval objects, not of rounded digits.
Thus$h_L=1008m$ and$269\nmid h_L$, as claimed.

For the final sector statement, the normal closure of$L$ is$M$ with group
$S_3\times C_2$, and$L$ is the fixed field of a reflection$s$.
Extension and norm of ideal classes have composites$2$ and$1+s$.
Since$2$ is invertible on the269-primary part and$\operatorname{Cl}(L)[269]=0$,
the$s$-invariant part of$\operatorname{Cl}(M)[269]$ is zero.
Every conjugate reflection acts as$-1$, and their products, the cubic
rotations, act as$1$.  Hence the nontrivial cubic sector is zero.
No assertion about the trivial cubic sector is made.
\end{proof}

\paragraph{Exact certificate and source scope.}
The full rational unit coefficients, monic polynomial identities,30 residue
rows and selected determinant99 are in
\texttt{saturation\_certificate.json}; the standalone independent verifier
is \texttt{verify\_saturation.py}.  PARI only proposed the units.
Its900-second and3600-second flag1 class-quotient certification jobs both
failed to certify;
their logs and receipts are preserved separately and are not evidence
for the conclusion above.
The analytic verifier \texttt{certify\_D30935\_analytic.py} and its exact
\texttt{D30935\_analytic\_interval.json} completed successfully in197seconds.
The primitive quadratic character table is independently compared at every
odd residue with the Jacobi symbol, and its finite Gauss sign is certified
in \texttt{verify\_quadratic\_gauss.py} and its output.

The analytic class-number formula and projected weighted-log regulator
normalization were freshly read on printed pages1 and4 of
A.~Sutherland, MIT18.785 Lecture19 (13 November2019),
\url{https://math.mit.edu/classes/18.785/2019fa/LectureNotes19.pdf}.
The classical theorem is retained, not reproved here.
The finite Fourier, modulus-four, Gaussian normalization, Mellin split,
elementary regulator and rational-isolation arguments were rederived
explicitly above.  Standard Poisson summation, Dirichlet's unit theorem
and the already retained Artin/ideal-character correspondence are used.
The python-flint0.9.0 inclusion arithmetic, outward radius convention and
\texttt{arb.expint} argument/order interface were freshly checked at
\url{https://python-flint.readthedocs.io/en/latest/arb.html}.
No full runtime implementation proof or Lean execution is claimed.

\paragraph{Geometric dictionary and remaining host gate.}
The cubic class-sector conclusion concerns unramified abelian covers of
$\operatorname{Spec}\mathcal O_M$, via the retained class-field dictionary
and the audited SGA1, ExposeV, sections4 and7 finite-cover correspondence.
Reflection invariants use exact prime-to269 extension and norm.
This certificate supplies the missing$D=30935$ arithmetic input of the
separate uncompensated e133 host audit.  Its use in a newly assembled
host, changed residual tuple, native support and original model still
requires those complete interfaces to be checked.  No automatic transfer
of the old e267 native ideal, and no H04.4 conclusion, follows here.
